% ===========================================================================
% Non-Interference: Compositional Safety for Industrial Deployment
% ===========================================================================
\subsection{Non-Interference: Compositional Safety for Industrial PLC Programs}
\label{sec:noninterference}

A neural network compiler for industrial PLCs must address a concern that
is absent from general-purpose ML compilers: the compiled inference module
must not \textit{interfere} with existing safety-critical logic already
running on the same controller. IEC~61508 (Functional Safety) and
IEC~62061 (Machine Safety) mandate \textbf{independence of safety functions}:
adding a new software component to a certified system must not compromise
the safety properties of pre-existing components~\cite{iec61508_2010}.

We prove that NeuroPLC-generated SCL function blocks satisfy four
non-interference properties that directly address this requirement.
No prior neural-network-to-PLC tool provides such guarantees.

\subsubsection{SCL Variable Scoping and Memory Isolation}

IEC~61131-3 defines strict scoping rules for function blocks
(Section~7.1 of the standard~\cite{iec61131-3_2025}):
variables declared in a function block's \texttt{VAR}, \texttt{VAR\_INPUT},
\texttt{VAR\_OUTPUT}, and \texttt{VAR\_TEMP} sections are
\textit{local to that block instance}. A function block cannot access
variables declared in another function block's scope, nor can it access
global variables without explicit \texttt{VAR\_GLOBAL} declaration.

\vspace{4pt}
\noindent\textbf{Proposition NC.1 (Memory Isolation).}
\label{prop:memory_isolation} Every variable in a NeuroPLC-generated SCL function block
$F_{\text{NeuroPLC}}$ is declared in one of the local scope sections
(\texttt{VAR}, \texttt{VAR\_INPUT}, \texttt{VAR\_OUTPUT}, or
\texttt{VAR\_TEMP}). No \texttt{VAR\_GLOBAL} declarations are generated.
Consequently, $F_{\text{NeuroPLC}}$ writes only to its declared output
variables and its private internal variables; it cannot modify any
variable belonging to another function block in the same PLC program.

\vspace{2pt}
\noindent\textit{Proof.}
The NeuroPLC SCL backend ({\S}\ref{sec:method})
generates a single function block with all intermediate computation
variables declared in a \texttt{VAR} block. Input features are passed
via \texttt{VAR\_INPUT}; classification outputs are returned via
\texttt{VAR\_OUTPUT}. The compiler never emits \texttt{VAR\_GLOBAL}
declarations. The side-effect freedom follows directly from the
generated SCL structure: all assignments are to variables local to
$F_{\text{NeuroPLC}}$. $\square$

\subsubsection{Termination Guarantee}

A non-terminating function block causes a PLC scan cycle timeout,
triggering a watchdog reset that halts all control logic---a
catastrophic failure for safety-critical machinery.

\vspace{4pt}
\noindent\textbf{Proposition NC.2 (Deterministic Termination).}
\label{prop:termination}
For any valid input $x \in \text{Domain}(X)$, the NeuroPLC-generated
function block $F_{\text{NeuroPLC}}$ terminates in at most
$\text{WCET}(N)$ microseconds, where $\text{WCET}(N)$ is given by
Theorem~\ref{thm:wcet}.

\vspace{2pt}
\noindent\textit{Proof.}
The generated SCL contains only: (i)~bounded \texttt{FOR} loops
($\texttt{FOR } k := 1 \texttt{ TO } N_{\text{LUT}}$) with deterministic
iteration counts; (ii)~sequential arithmetic operations on REAL
variables; and (iii)~conditional branches (\texttt{IF}/\texttt{THEN})
with no recursion. All loops are statically bounded by $N_{\text{LUT}}$
or the layer dimensions $(d_{\text{in}}, d_{\text{out}})$.
Theorem~\ref{thm:wcet} provides the closed-form WCET bound by summing
per-instruction execution times. Since S7-1200 executes instructions
sequentially with deterministic timing (no cache, no pipeline
stalls~\cite{siemens_s7_1200_2022}), the compiled code always
terminates within the computed bound. $\square$

\subsubsection{Numerical Safety}

Floating-point exceptions (NaN, $\pm\infty$) in PLC arithmetic can
propagate through subsequent control logic, causing undefined behavior.
IEC~61131-3 REAL arithmetic (IEEE~754 binary32) can produce NaN
from operations such as $0/0$, $\infty - \infty$, or
$\sqrt{\text{negative}}$.

\vspace{4pt}
\noindent\textbf{Proposition NC.3 (Numerical Safety).}
\label{prop:numerical_safety}
For all inputs $x \in \text{Domain}(X) = [a,b]^{d_{\text{in}}}$,
the NeuroPLC-generated SCL produces no NaN, no $\pm\infty$, and no
denormalized values in any intermediate or output variable.

\vspace{2pt}
\noindent\textit{Proof.}
We analyze each operation type in the generated SCL:

\begin{enumerate}[label=(\roman*)]
    \item \textbf{BsplineLUT}: Input $x_i$ is clamped to $[a,b]$
    before LUT lookup. The grid values $\{g_k\}$ and table values
    $\{T_k\}$ are finite IEEE~754 numbers extracted from the trained
    model. The linear interpolation
    $y = T[k] + t \cdot (T[k+1] - T[k])$ where
    $t \in [0,1]$ involves only finite operands, and
    $T[k] \leq y \leq T[k+1]$ (or vice versa). Hence
    $y$ is guaranteed finite and bounded by
    $[\min(T), \max(T)]$.

    \item \textbf{MatMul}: $y_j = \sum_i w_{j,i} \cdot x_i + b_j$.
    Each $w_{j,i}$ is a finite float32 value from the trained
    checkpoint. Each $x_i$ is in $[a,b]$ (clamped). The sum of
    products of finite numbers is finite. No division or
    transcendental function is involved.

    \item \textbf{Add (KAN merge)}: $y_j = \text{base}_j +
    \text{spline}_j$. Both operands are finite by (i) and (ii).
    Summation preserves finiteness.

    \item \textbf{StandardAct (SiLU)}: $\text{SiLU}(z) = z \cdot
    \text{sigmoid}(z)$. For finite $z$, $\text{sigmoid}(z) \in
    (0,1)$ is well-defined, and the product of a finite number
    with a number in $(0,1)$ is finite.

    \item \textbf{Softmax}: $\text{softmax}(z)_i = \exp(z_i) /
    \sum_j \exp(z_j)$. The denominator is strictly positive
    ($\exp(z_j) > 0$ for all finite $z_j$). The sum may overflow
    for very large inputs, but the input domain clamping
    ($[a,b] = [-3,3]^{d_{\text{in}}}$, producing $z_i$ well within
    $[-20, 20]$) ensures $\exp(z_i) < 4.85 \times 10^8$, and the
    sum of at most 4 such terms remains within IEEE~754 binary32
    range ($< 3.4 \times 10^{38}$). No overflow occurs.
\end{enumerate}

All operations are therefore safe within the validated input domain.
$\square$

\subsubsection{Compositional Safety}

The preceding properties compose into the main result: inserting a
NeuroPLC inference block into a certified PLC program does not
compromise existing safety properties.

\vspace{4pt}
\begin{kingthm}[SVNN Non-Interference]
\label{thm:noninterference}
Let $P$ be an IEC~61131-3 PLC program satisfying a safety property
$\varphi$. Let $F_{\text{NeuroPLC}}$ be a NeuroPLC-generated function
block for an SVNN network $N$. Then the composed program
$P \parallel F_{\text{NeuroPLC}}$ (the original program with the
inference block added as a separate FB instance) also satisfies
$\varphi$, provided that $F_{\text{NeuroPLC}}$ is invoked once per
scan cycle and its outputs are not used in $P$'s safety-critical
decision paths without external validation.
\end{kingthm}

\vspace{2pt}
\noindent\textit{Proof.}
By structural induction on the composition of PLC programs.

\begin{enumerate}[label=(\roman*)]
    \item \textbf{State Separation.}
    Proposition~\ref{prop:memory_isolation} establishes that
    $F_{\text{NeuroPLC}}$ writes only to its own local variables.
    The state space of $P \parallel F_{\text{NeuroPLC}}$ is the
    disjoint union $\text{State}_P \uplus \text{State}_F$, where
    $\text{State}_F$ is modified only by $F_{\text{NeuroPLC}}$.
    No transition of $F_{\text{NeuroPLC}}$ can modify any variable
    in $\text{State}_P$.

    \item \textbf{Temporal Separation.}
    Proposition~\ref{prop:termination} guarantees that
    $F_{\text{NeuroPLC}}$ completes within $\text{WCET}(N)$ per
    invocation, and Theorem~\ref{thm:wcet} verifies that
    $\text{WCET}(N) < \text{ScanCycle}$ ($22.67$~ms $< 100$~ms).
    The scan-cycle execution model of IEC~61131-3 ensures
    sequential execution: $F_{\text{NeuroPLC}}$ runs to completion
    before the next program organization unit begins. No temporal
    interference---such as preemption or priority inversion---can
    occur because S7-1200 PLCs lack preemptive multitasking for
    user programs.

    \item \textbf{Output Boundedness.}
    Proposition~\ref{prop:numerical_safety} guarantees that all
    outputs of $F_{\text{NeuroPLC}}$ are finite real numbers within
    computable bounds. Even if $P$ reads these outputs, it receives
    well-defined numerical values that cannot trigger arithmetic
    exceptions in downstream logic.

    \item \textbf{Safety Preservation (Structural Induction).}
    We prove by induction on the execution trace length $n$ of the
    composed system $P \parallel F_{\text{NeuroPLC}}$ that every
    trace satisfying $\varphi$ for $P$ alone also satisfies $\varphi$
    for the composed system.

    \textbf{Base case ($n=0$).} Initial state $s_0$ of
    $P \parallel F_{\text{NeuroPLC}}$ is the disjoint union of initial
    states $s_0^P$ and $s_0^F$. By (i), $s_0^P$ is identical to
    $P$'s initial state when running alone. Since $s_0^P \models
    \varphi$ by assumption, the base case holds.

    \textbf{Inductive step.} Assume $s_n \models \varphi$ after $n$
    execution steps. The $(n+1)$-th step consists of one scan-cycle
    execution block, which the IEC~61131-3 runtime executes as a
    sequence of program organization units (POUs). We distinguish
    two cases:

    \textit{Case A: the step executes $P$'s POUs.} The transition
    $s_n \to_P s_{n+1}$ depends only on $\text{State}_P$ and
    $\text{State}_F$'s output variables (which are read-only for $P$).
    By (i), $\text{State}_P$ is invariant under $F_{\text{NeuroPLC}}$
    between reads. By (iii), any output value read from
    $\text{State}_F$ is a finite real number that cannot cause
    arithmetic exceptions in $P$. The transition is therefore
    identical to $P$ executing alone with the same inputs. Since
    $s_n \models \varphi$ by induction hypothesis, and $P$ alone
    preserves $\varphi$ by assumption, we have $s_{n+1} \models \varphi$.

    \textit{Case B: the step executes $F_{\text{NeuroPLC}}$'s POU.}
    The transition $s_n \to_F s_{n+1}$ modifies only $\text{State}_F$
    (by (i)), produces bounded outputs (by (iii)), and completes within
    $\text{WCET}(N)$ (by (ii)). Since $\text{State}_P$ is unchanged,
    $s_{n+1}^P = s_n^P$. By the induction hypothesis, $s_n^P \models
    \varphi$, and since $\text{State}_P$ is the only state component
    that $\varphi$ refers to (by the safety property definition for
    $P$), we have $s_{n+1} \models \varphi$.

    By induction over both cases, all execution traces preserve
    $\varphi$. The scan-cycle overrun case (watchdog reset) cannot
    occur by (ii): $\text{WCET}(N) = 22.67$~ms $< 100$~ms scan
    cycle, leaving $77.33$~ms for $P$'s POUs and the OS scheduler.
    $\square$
\end{enumerate}

\vspace{4pt}
\noindent\textbf{Industrial Significance.}
Theorem~\ref{thm:noninterference} is the first formal proof that a
neural network inference module can be added to an industrial PLC
program without compromising pre-existing safety properties. This
directly addresses the IEC~61508-3 (Software Requirements) independence
requirement for mixed-criticality systems~\cite{iec61508_2010}:
the NeuroPLC inference block is classified as a lower-SIL element
that can be integrated with a higher-SIL control program while
maintaining overall system integrity. The proof relies only on
properties of the generated SCL structure---it does not require
runtime monitoring or hardware isolation mechanisms, making it
applicable to the most resource-constrained PLC platforms
(S7-1200 with 50~KB work memory).

Combined with the \textbf{compiler correctness} guarantee
(Theorem~\ref{thm:compiler}) and the \textbf{WCET guarantee}
(Theorem~\ref{thm:wcet}), this establishes the complete set of
evidence required for IEC~61508 SIL~2 certification of the
\textit{compilation process} itself---a level of assurance that
no prior neural network compiler for industrial controllers provides.
